% ---- pg-317 （书页 305） ----

\begin{mvframe}[显示板 7.2(续)]
\begin{center}
\begin{tabular}{lllll}
\toprule
Parameter & Estimate & T for H0: Parameter $=0$ & Pr $>|T|$ & Std Error of Estimate\\
\midrule
INTERCEPT & $-1.000\ 000\ 000$ & $-1.12$ & $0.345\ 0$ & $0.894\ 427\ 19$\\
Z1 & $1.000\ 000\ 000$ & $2.74$ & $0.071\ 4$ & $0.365\ 148\ 37$\\
\bottomrule
\end{tabular}

\medskip
\begin{tabular}{c}
E = Error SS \& CP Matrix
\end{tabular}

\medskip
\begin{tabular}{ccc}
 & Y1 & Y2\\
Y1 & 6 & $-2$\\
Y2 & $-2$ & 4\\
\end{tabular}

\medskip
\begin{tabular}{c}
Manova Test Criteria and Exact F Statistics for\\
the Hypothesis of no Overall Z1 Effect\\
H = Type III SS\&CP Matrix for Z1 \quad E = Error SS\&CP Matrix\\
S = 1 \quad M = 0 \quad N = 0
\end{tabular}

\medskip
\begin{tabular}{lrrrrr}
\toprule
Statistic & Value & F & Num DF & Den DF & Pr $>$ F\\
\midrule
Wilks' Lambda & 0.062 500 00 & 15.000 0 & 2 & 2 & 0.062 5\\
Pillai's Trace & 0.937 500 00 & 15.000 0 & 2 & 2 & 0.062 5\\
Hotelling-Lawley Trace & 15.000 000 00 & 15.000 0 & 2 & 2 & 0.062 5\\
Roy's Greatest Root & 15.000 000 00 & 15.000 0 & 2 & 2 & 0.062 5\\
\bottomrule
\end{tabular}
\end{center}
\end{mvframe}

\mvthm{结论 7.9}{对模型(7-23)及满秩 $(\mathbf{Z})=r+1<n$, 最小二乘估计量
$\hat{\boldsymbol{\beta}}=\left[\begin{array}{c:c:c:c}\hat{\boldsymbol{\beta}}_{(1)}&\hat{\boldsymbol{\beta}}_{(2)}&\cdots&\hat{\boldsymbol{\beta}}_{(m)}\end{array}\right]$
满足}
\[
E(\hat{\boldsymbol{\beta}}_{(i)})=\boldsymbol{\beta}_{(i)},\qquad E(\hat{\boldsymbol{\beta}})=\boldsymbol{\beta}
\]
且
\[
\operatorname{Cov}(\hat{\boldsymbol{\beta}}_{(i)},\hat{\boldsymbol{\beta}}_{(k)})=\sigma_{ik}(\mathbf{Z}'\mathbf{Z})^{-1},\qquad i,k=1,2,\cdots,m
\]
残差
$\hat{\boldsymbol{\epsilon}}=\left[\begin{array}{c:c:c:c}\hat{\boldsymbol{\epsilon}}_{(1)}&\hat{\boldsymbol{\epsilon}}_{(2)}&\cdots&\hat{\boldsymbol{\epsilon}}_{(m)}\end{array}\right]=\mathbf{Y}-\mathbf{Z}\hat{\boldsymbol{\beta}}$,
满足 $E(\hat{\boldsymbol{\epsilon}}_{(i)})=\mathbf{0}$ 且 $E(\hat{\boldsymbol{\epsilon}}_{(i)}'\hat{\boldsymbol{\epsilon}}_{(k)})=(n-r-1)\sigma_{ik}$, 故
\[
E(\hat{\boldsymbol{\epsilon}})=\mathbf{0}\qquad E\left(\frac{1}{n-r-1}\hat{\boldsymbol{\epsilon}}'\hat{\boldsymbol{\epsilon}}\right)=\boldsymbol{\Sigma}
\]
此外, $\hat{\boldsymbol{\epsilon}}$ 和 $\hat{\boldsymbol{\beta}}$ 不相关.

\medskip
\noindent\textbf{证明}\quad 第 $i$ 个响应服从多重回归模型
\[
\mathbf{Y}_{(i)}=\mathbf{Z}\boldsymbol{\beta}_{(i)}+\boldsymbol{\epsilon}_{(i)},\quad
E(\boldsymbol{\epsilon}_{(i)})=\mathbf{0},\quad
E(\boldsymbol{\epsilon}_{(i)}\boldsymbol{\epsilon}_{(i)}')=\sigma_{ii}\mathbf{I}
\]
又, 如同\hyperref[eq:7-24]{式(7-24)}中那样有
\phantomsection\label{eq:7-33}
\begin{equation}
\hat{\boldsymbol{\beta}}_{(i)}-\boldsymbol{\beta}_{(i)}
=(\mathbf{Z}'\mathbf{Z})^{-1}\mathbf{Z}'\mathbf{Y}_{(i)}-\boldsymbol{\beta}_{(i)}
=(\mathbf{Z}'\mathbf{Z})^{-1}\mathbf{Z}'\boldsymbol{\epsilon}_{(i)}
\tag{7-33}
\end{equation}
及
\[
\hat{\boldsymbol{\epsilon}}_{(i)}=\mathbf{Y}_{(i)}-\hat{\mathbf{Y}}_{(i)}
=[\mathbf{I}-\mathbf{Z}(\mathbf{Z}'\mathbf{Z})^{-1}\mathbf{Z}']\mathbf{Y}_{(i)}
=[\mathbf{I}-\mathbf{Z}(\mathbf{Z}'\mathbf{Z})^{-1}\mathbf{Z}']\boldsymbol{\epsilon}_{(i)}
\]
故 $E(\hat{\boldsymbol{\beta}}_{(i)})=\boldsymbol{\beta}_{(i)}$ 且 $E(\hat{\boldsymbol{\epsilon}}_{(i)})=\mathbf{0}$. 接下来, 有
\[
\begin{aligned}
\operatorname{Cov}(\hat{\boldsymbol{\beta}}_{(i)},\hat{\boldsymbol{\beta}}_{(k)})
&=E(\hat{\boldsymbol{\beta}}_{(i)}-\boldsymbol{\beta}_{(i)})(\hat{\boldsymbol{\beta}}_{(k)}-\boldsymbol{\beta}_{(k)})'\\
&=(\mathbf{Z}'\mathbf{Z})^{-1}\mathbf{Z}'E(\boldsymbol{\epsilon}_{(i)}\boldsymbol{\epsilon}_{(k)}')\mathbf{Z}(\mathbf{Z}'\mathbf{Z})^{-1}
=\sigma_{ik}(\mathbf{Z}'\mathbf{Z})^{-1}
\end{aligned}
\]
用结论 4.9 及结论 7.2 的证明, 对任意随机向量 $\mathbf{U}$ 和常数阵 $\mathbf{A}$, 我们有
$E[\mathbf{U}'\mathbf{A}\mathbf{U}]=$
